% ===========================================================================
\section{Conclusion}
\label{sec:conclusion}
% ===========================================================================

On common FY2023--24 production support, the legacy and refined methods give
\SI{3.317}{} and \SI{4.485}{\mega\tonne} \COe{}. Expanding to the public-data 2019 segment
ledger raises the conservative (default-practice) configuration to
\SI{23.499}{\mega\tonne}---high in its own 9.601--24.694 joint envelope, whose top is the
same practice setting with fertilizer feedstock inside the post-meter leak base. The
reference should be read as an
upper-practice conservative configuration and never as an empirically centred
best estimate: no probability attaches to any point in the envelope, and until
Pakistani pipe-material, pressure, LDAR or post-meter measurements exist, the
honest statement is the envelope with its selection rule, with the joint floor
priced beside it. The principal change is
therefore supply-chain coverage rather than the production-factor revision alone.

Joint structural choices span \SIrange{9.601}{24.694}{\mega\tonne}. Distribution technology
and transferred post-meter industrial leakage determine most of this range and
\SI{90.4}{\percent} of the reference-configuration rank sensitivity. Utility balances show that
the network estimates are physically feasible, while distinguishing sharply between higher- and
lower-practice interpretations.

Both adjacent-fiscal-year bracket endpoints are more than five times Pakistan's 2022 CRT
on a common GWP basis;
the difference traces to activity, factor and boundary choices rather than arithmetic. Satellite
observations confirm broad upstream source geography but show that production weakly predicts
detected intensity. National pipe characteristics, post-meter measurements and plant gas
disposition are therefore the highest-value next data. The implemented ledger can incorporate
those records without another methodological redesign.
